\scanpage{127}
\mspage{114}
Quant à la flèche $\widetilde\beta_0\colon\widetilde B\unsure{\text{---}|\longrightarrow}Y_0$, elle est
donnée par les conditions de commutativité
\[
 \widetilde\beta_0r=\unsure{p_0|\beta_0}\alpha_0,\qquad
 \widetilde\beta_0s=\lambda_Y\beta_{I'},\tag*{(6.3.11)}
\]
ce qui est \inserted{encore} licite, puisque
\[
 (p_0\alpha_0)\lambda_A=(\lambda_Y\beta_{I'})i'_{I'}
\]
en vertu de la commutativité du
cube (6.3.6). Pour établir la commutativité
totale du diagramme (6.3.9)\unsure{,}
(traits pleins, et \deleted{traits}\unsure{,} discontinus
pour $h'_{I'}$ \deleted{et}, $\widetilde h_0$, et $\widetilde\beta_0$ -- à l'exclusion
de $h_0$ pas \unsure{encore|encor} construit), il faut
encore vérifier \inserted{d'une part \deleted{les}} les relations
\[
 \left\{\begin{aligned}
 \widetilde h_0r&=\alpha_0,\qquad \underline{p_0\widetilde h_0=\widetilde\beta_0},\\
 \lambda_Xh'_{I'}&=\widetilde h_0s
 \end{aligned}\right.\qquad
 \deleted{\text{(déjà connue, cf }(6.3.10)\text{)}}
 \tag*{(6.3.12)}
\]
% slanted
\margin{NB la première et la troisième de ces relations sont déjà connues, cf (6.3.10)}
établissant la commutativité du
cube supérieur avec le carré diagonal
(diagramme du type (6.3.6), la flèche
pointillée $h_0$ y incluse), d'autre
part les relations
\[
 \left\{\begin{aligned}
 \widetilde i r&=i_0\quad(\text{et }\widetilde i s=\lambda_B\ \text{pour mémoire})\\
 \beta_0\widetilde i&=\widetilde\beta_0,
 \end{aligned}\right.\tag*{(6.3.13)}
\]
où la première ligne est déjà connue
par définition de $\widetilde B$ et de $\widetilde i$. La deuxième
% archive sheet number: 126

